% In Catilinam II (in-catilinam-2) --- English translation
% Translated by Claude (Anthropic) from the Latin (Perseus).
% Style guide: STYLE.md.

\ciceroSection{1} At long last, citizens, Lucius Catiline --- raging with audacity, panting crime, contriving wickedly the ruin of his country, threatening you and this city with sword and fire --- we have either cast out, or sent forth, or, as he was leaving of his own motion, accompanied with parting words. He has gone, he has departed, he has escaped, he has burst out. No ruin will any longer be prepared by that monster and prodigy of a man within the walls against the very walls. And this man, the one undisputed leader of this domestic war, we have without contention defeated. For now that dagger of his shall not turn between our ribs; not in the Field of Mars, not in the Forum, not in the Senate-house, not finally within our own walls shall we tremble. He was moved from his ground when he was driven from the city. Openly, now, with no one in the way, we shall wage a regular war against an enemy. Without doubt we have undone the man and conquered him magnificently, when we have flung him out of his hidden ambush into open brigandage.

\ciceroSection{2} But because he did not, as he wished, carry forth his blade bloody, because he came forth while we yet lived, because we wrenched the steel from his hands, because he left the citizens unharmed and the city standing --- with how great a grief, then, do you suppose he is afflicted and brought low? He lies prostrate now, citizens, and feels himself struck down and abject; and he turns his eyes again and again, no doubt, toward this city which he mourns has been snatched from his jaws --- which seems to me to rejoice that it has vomited out so foul a plague and cast it forth.

\ciceroSection{3} And if there is any such man as all should have been, who in this very thing in which my speech exults and triumphs vehemently accuses me --- because I did not arrest so capital an enemy rather than let him go --- the fault is not mine, citizens, but the times'. Lucius Catiline ought long ago to have been put to death and visited with the heaviest punishment, and this our ancestral custom and the severity of this empire and the commonwealth demanded of me. But how many do you suppose there were who would not believe what I reported, how many who from folly thought it not so, how many who would even defend him, how many who from wickedness favoured him? And if I had judged that with him removed every danger could be driven from you, I should long since have removed Lucius Catiline at the risk not only of odium against me but even of my life.

\ciceroSection{4} But when I saw that, with not even all of you yet then convinced, if I had visited him as he deserved with death, I should be unable, crushed by odium, to pursue his confederates, I led the matter to this point: that you might be able to fight openly when you saw the enemy openly. Just how vehemently this enemy must, in my judgment, be feared abroad, you may understand from this --- that I take it amiss even that he has gone out of the city with too small a following. Would that he had taken with him all his forces! Tongilius he took out for me, with whom he had begun a love-affair when the boy still wore the praetexta; Publicius and Minucius, whose debt incurred in a tavern could bring no commotion to the commonwealth. He left behind --- what men, with how much debt, how strong, how noble!

\ciceroSection{5} And so I despise that army of his, set against the Gallic legions and against this levy which Quintus Metellus has held in the Picene and Gallic country, and against these forces which are being daily prepared by us --- gathered as it is from desperate old men, from rustic luxury, from country bankrupts, from those who have preferred to default on their bail-bonds rather than that army of his. To these, if I show not the line of our army but even the praetor's edict, they will collapse. These men whom I see flitting in the Forum, standing at the Senate-house, even coming into the Senate, who shine with ointments, who glitter in purple --- I would rather he had taken his own soldiers with him: who, if they remain here, remember it is not so much that army that we must fear as these men who have deserted that army. And these men are the more to be feared because they perceive that I know what they are planning, and yet they are not at all moved.

\ciceroSection{6} I see to whom Apulia has been allotted, who has Etruria, who the Picene country, who the Gallic, who has demanded for himself this work of urban plot, of slaughter and of fires. They perceive that all the plans of last night have been brought to me; I laid them open in the Senate yesterday; Catiline himself was struck with terror and fled. These men --- what are they waiting for? They greatly err if they hope that former mildness of mine will be perpetual. What I have waited for I have already attained --- that you all should see openly that a conspiracy has been made against the commonwealth; unless of course there is anyone who thinks that men like Catiline do not feel with Catiline. There is no place now for mildness; the matter itself demands severity. One concession even now I shall make: let them go, let them set out, let them not allow Catiline to waste away in misery from longing for them. I shall point out the road: he has set out by the Aurelian Way; if they wish to hurry, they will catch him by evening.

\ciceroSection{7} O fortunate commonwealth, if it shall have cast out this bilge of the city! With Catiline alone drawn off, by Hercules, the commonwealth seems to me lightened and refreshed. For what evil or crime can be imagined or conceived which he has not contrived? What poisoner in all Italy, what gladiator, what brigand, what cut-throat, what parricide, what forger of wills, what swindler, what glutton, what spendthrift, what adulterer, what woman of ill fame, what corrupter of youth, what corrupted, what abandoned man can be found, who would not confess he has lived with Catiline on terms of the closest intimacy? What slaughter through these years has been done without him, what unspeakable lust not through him?

\ciceroSection{8} What enticement to youth was ever in any man so great as in him? who himself loved some most basely, served the love of others most scandalously, promised some the fruit of their lusts, others the death of their parents --- not only urging it but even helping to it. Now indeed how suddenly he had collected, not from the city only but even from the country, a vast number of lost men! No one weighed down by debt, not only at Rome but in any corner of all Italy, has there been whom he has not enrolled into this incredible league of crime.

\ciceroSection{9} And that you may perceive his diverse pursuits in their unlike methods --- there is no one in the gladiators' school somewhat readier for crime than the rest who does not confess himself an intimate of Catiline; no one on the stage more frivolous and worthless who does not recall having been almost a comrade of his. And the same man, accustomed by the practice of debaucheries and crimes to the bearing of cold and hunger and thirst and sleeplessness, was praised by these men as ``hardy'' --- when he was using up the supports of industry and the instruments of virtue on his lust and audacity.

\ciceroSection{10} If, indeed, his comrades shall follow him, if these scandalous herds of desperate men shall go out of the city --- O happy us, O fortunate commonwealth, O distinguished praise of my consulship! For these men's lusts are not now middling, their audacity not human and bearable; they think of nothing but slaughter, but fires, but plunder. They have squandered their patrimonies, they have mortgaged their fortunes; their property long since, their credit lately, has begun to fail them. Yet that same lust which there was in their abundance remains. If in their drink and dice they sought only feasts and harlots, they would indeed be men to be despaired of, but yet to be borne. But who could bear this --- that idle men should plot against the bravest, the most foolish against the most prudent, the drunken against the sober, the sleeping against the waking? who, reclining at my banquets, embracing shameless women, languid with wine, stuffed with food, garlanded with wreaths, smeared with ointments, weakened by debauches, belch out in their conversations the slaughter of loyal men and the burning of the city.

\ciceroSection{11} Whom I am confident that some doom hangs over, and that the punishment owed long since to wickedness, worthlessness, crime, and lust is now plainly upon them or certainly draws near. Whom if my consulship, since it cannot heal them, shall have removed --- it shall have prolonged for the commonwealth not some short time but many ages. For there is no nation we should fear, no king who can make war upon the Roman people. All abroad has been pacified by one man's courage, by land and sea: there remains domestic war, the snares are within, within is the danger enclosed, within is the enemy. With luxury, with madness, with crime we must contend. To this war I declare myself the leader, citizens; I take up the enmities of lost men. What can be cured I shall cure by any method; what must be cut away I shall not allow to remain to the ruin of the city. So let them either go out, or be quiet, or, if they remain in the city and in the same mind, let them expect what they deserve.

\ciceroSection{12} But there are still those who say, citizens, that Catiline was cast out by me. If I could attain it by a word, I would cast out those very men who say this. For evidently the timid or even very modest man could not bear the consul's voice; as soon as he was bidden to go into exile, he obeyed. Why, even yesterday, when I had nearly been killed at home, I called the Senate to the temple of Jupiter Stator and laid the whole matter before the senators. When Catiline came there, what senator addressed him, who saluted him, who at last so much as looked at him as a lost citizen and not rather as a most dangerous enemy? --- nay, even the leaders of that order left bare and empty that part of the benches to which he had drawn near.

\ciceroSection{13} Then I, that vehement consul who casts out citizens into exile by a word, asked Catiline whether he had been at the night-meeting at Marcus Laeca's or no. When that most audacious man, convicted by his conscience, was at first silent, I laid open the rest: what he had done that night, where he had been, what he had appointed for the next, in what way the plan of the entire war had been laid out for him --- I expounded all. When he hesitated, when he was caught, I asked what he hesitated to set out where he had long been preparing, when I knew that arms, axes, fasces, trumpets, military standards, and that silver eagle for which he had even set up a shrine in his house --- had been sent ahead.

\ciceroSection{14} Was I casting into exile the man whom I had seen had already entered upon war? For, I suppose, that centurion Manlius who pitched camp in the Faesulan country has on his own authority declared war on the Roman people, and that camp is not now expecting Catiline as its leader, and the man cast into exile is taking himself, as they say, to Massilia, not to that camp. O wretched condition not only of administering but even of preserving the commonwealth! Now if Lucius Catiline, hemmed in and weakened by my counsels, my labours, my dangers, shall suddenly take fright, change his mind, desert his men, abandon the design of war, and from this course of crime and war turn his way to flight and exile, he will not be said to have been stripped of the arms of his audacity by me, not stupefied and terrified by my diligence, not driven from his hope and his attempt; he will be said to have been cast into exile, uncondemned and innocent, by force and threats of the consul: and there will be those who, if he does this, would have him reckoned not wicked but wretched, and me reckoned not the most diligent consul but the cruellest tyrant!

\ciceroSection{15} It is worth as much to me, citizens, to undergo the storm of this false and unjust odium, provided only that the danger of this dreadful and wicked war be turned aside from you. By all means let it be said that he was cast out by me, provided only that he go into exile. But, believe me, he will not go. Never, citizens, shall I pray to the immortal gods, for the relief of my odium, that you should hear of Lucius Catiline leading an army of enemies and flitting about in arms; yet within three days you will hear it. And I am much more afraid of this --- that it may yet be invidious against me that I let him out rather than cast him out. But when there are men who say that, since he set out, he was cast out --- the same, if he had been killed, what would they say?

\ciceroSection{16} And those who keep saying that Catiline is on his way to Massilia complain of it less than they fear it. There is no one of them so merciful as not to prefer that he go to Manlius rather than to the Massilians. As for him, by Hercules, even if he had never before thought of what he is now doing, yet he would prefer being killed in his brigandage to living in exile. Now, indeed, since nothing has so far befallen him beyond his own will and design, except that he set out from Rome while we yet lived, let us rather hope that he go into exile than complain.

\ciceroSection{17} But why do we speak so long of one enemy --- and of one enemy who now confesses he is an enemy, and whom, since (what I always wished) a wall stands between us, I do not fear? --- and say nothing of those who dissemble, who remain at Rome, who are with us? whom indeed I, if it can in any way be done, am eager not so much to punish as to heal --- to themselves, to placate to the commonwealth; nor do I see why this could not be done, if they are willing now to hear me. For I shall set out for you, citizens, of what kinds of men these forces are gathered; then to each I shall apply the medicine of my counsel and my speech, if I can.

\ciceroSection{18} One kind is of those who, with a great debt, have yet greater holdings, and from love of which they cannot in any way bring themselves to come loose. The look of these men is most respectable --- for they are well-off --- but their will and motive are most shameless. You are equipped and rich in lands, in buildings, in silver, in slaves, in everything; and yet you hesitate to take from your holding to add to your credit? For what are you waiting? War? What then? --- in the laying-waste of all things do you suppose your own holdings will be sacrosanct? Or do you wait for new account-books? They err who expect such from Catiline: by my kindness new account-books are set up --- but auctioneer's books; for these men who have holdings can in no other way be safe. If they had been willing to do it earlier, and not (which is most foolish) to fight the interest with the produce of their estates, we should be using them as both wealthier and better citizens. But these men I think least to be feared, because they can either be drawn off from their resolve, or, if they remain in it, seem to me likelier to make vows against the commonwealth than to bear arms.

\ciceroSection{19} The second kind is of those who, although weighed down by debt, yet expect domination, wish to seize on power, think they can attain in a disturbed commonwealth the offices they despair of in a quiet one. To these the same precept seems to be given as to all the rest --- that they despair of attaining what they attempt: first of all, that I myself keep watch, am present, am providing for the commonwealth; then, that there are great spirits in good men, great concord of the orders, the greatest multitude, great forces of soldiers besides; that the immortal gods, finally, will bring to this unconquered people, this most distinguished empire, this most beautiful city, present help against so great a force of crime. But if they were now to attain that which in the highest madness they desire --- in the ash of the city and in the blood of the citizens, which they have, in their criminal and unspeakable mind, longed for, do they hope to be consuls or dictators or even kings? Do they not see that they desire what, if they attained it, they must of necessity yield up to some runaway slave or gladiator?

\ciceroSection{20} The third kind is of men now affected by age, but yet sturdy by exercise; of which kind is that Manlius, whose place Catiline is now taking. These are men from those colonies which Sulla established, which I see as wholes to be of the best citizens and bravest men; but yet they are colonists who threw themselves about, in moneys unhoped for and sudden, more lavishly and more insolently than was decent. While they build like men of fortune, while they delight in choice estates, large households, fitted-out banquets, they fell into so great a debt that, if they wished to be safe, Sulla would have to be raised for them from the dead. They have even driven some country-folk, slender and needy, into that same hope of the old plunder. Both of these I place in the same kind of plunderers and despoilers; but I warn them: let them cease to rage and to think of proscriptions and dictatorships. For so great a sorrow of those times has been seared into the state that not only men but not even, it seems to me, beasts will now bear them.

\ciceroSection{21} The fourth kind is mixed indeed, and motley, and turbulent: those who have long been weighed down, who never come to the surface, who partly through laziness, partly by mismanaging their business, partly by their expenses, totter in old debt, who, worn out by bail-bonds, by trials, by the proscription of their goods, are very many, and are said to be betaking themselves both from the city and from the country to that camp. These I think are not so much keen soldiers as slow defaulters. Let these men, the sooner the better, if they cannot stand, fall --- but in such a way that not only the state but not even their nearest neighbours feel it. For this I do not understand: why, if they cannot live honourably, they would wish to perish basely; or why they think they shall perish with less pain in company with many than if they perish alone.
